\section{Construction of Conversational Threads in Forums Data} 
\label{sec:how-to-thread}


Content comes from Reddit's data API in two separate but linked forms: \textit{submissions} and \textit{comments}. \textit{Submissions} are either "link posts" to external content (e.g. news articles, blogs, or even multimedia content) or "self posts" (submissions written by the poster meant to initiate a discussion thread on a topic). \textit{Comments} are user replies to either the initiating post (top level comments) or to another user's comment. Posts, top-level comments, and replies to comments form a nested conversational thread with a submission post at it's root and comments branching out into multiple possible dialogue trees. 





The tree-like structure of Reddit threads allows for multiple possible data formats depending on how the various components of a thread are combined. 
We investigate three formats for their potential as LM pretraining data:
\begin{itemize}[leftmargin=1em]
    \item \textbf{Atomic content}. This simple format treats all comments and submissions as independent documents without any structure or connection to the thread they appear in.
    \item \textbf{Partial threads}. This format assembles comments from the same thread into a structured, multi-round dialogue between users. Submissions are left as separate documents. Assembled dialogues are limited to a maximum parent depth, and the resulting documents are only snippets of a their originating thread (which are spread across several documents).
    \item \textbf{Full threads}. This complex format combines a given submission and all of its child comments into a single document encompassing an entire thread. Code-like indentation is used to indicate the depth of a comment in the thread's hierarchy.
\end{itemize}


We experimentally evaluated these strategies for assembling documents in \autoref{fig:abl_reddit}. We found that, for language modeling purposes, treating comments and submissions as atomic units leads to better downstream performance compared to partial and full threads. 
We hypothesize that the more complex formatting required to handle dialogues might introduce undesirable content for language modeling, such as short and repeated comments. 
We leave the study of better formatting for forum content for language modeling to future work.
